\chapter{Quasilinear Maximum Principles}

Consider the quasilinear operator
\begin{equation}
Qu=a^{ij}(x,u,Du)D_{ij}u+b(x,u,Du),
\qquad a^{ij}=a^{ji}.
\tag{10.1}
\end{equation}
Two such operators are equivalent if one is a positive multiple of the other.
On $\mathcal U\subset\Omega\times\R\times\R^n$, ellipticity means
\begin{equation}
0<\lambda(x,z,p)|\xi|^2
\le a^{ij}(x,z,p)\xi_i\xi_j
\le\Lambda(x,z,p)|\xi|^2
\tag{10.2}
\end{equation}
for every $\xi\ne0$ and $(x,z,p)\in\mathcal U$; uniform ellipticity means
that $\Lambda/\lambda$ is bounded there.  Ellipticity with respect to a
function $u$ means positivity along $(x,u(x),Du(x))$.  Set
\begin{equation}
\mathcal E(x,z,p)=a^{ij}(x,z,p)p_ip_j.
\tag{10.3}
\end{equation}
Then ellipticity implies
\begin{equation}
0<\lambda(x,z,p)|p|^2\le\mathcal E(x,z,p)
\le\Lambda(x,z,p)|p|^2
\tag{10.4}
\end{equation}
when $p\ne0$.

The operator is in divergence form if
\begin{equation}
Qu=\operatorname{div}A(x,u,Du)+B(x,u,Du),
\tag{10.5}
\end{equation}
so that
$a^{ij}=\tfrac12(D_{p_i}A^j+D_{p_j}A^i)$.  It is variational when, for
some integrand $F$,
\begin{equation}
A^i=D_{p_i}F,\qquad B=-D_zF.
\tag{10.6}
\end{equation}
In this case ellipticity is strict convexity of $F$ in the gradient variable.

The following model operators fix notation used below.  For $\alpha\ge1$,
\[
Qu=\Delta u+(\alpha-2)\frac{D_iuD_ju}{1+|Du|^2}D_{ij}u
\]
is equivalent to
$(1+|Du|^2)^{1-\alpha/2}\operatorname{div}
((1+|Du|^2)^{\alpha/2-1}Du)$ and is variational for
$\int(1+|Du|^2)^{\alpha/2}$.  Its cases $\alpha=2$ and $\alpha=1$ are,
respectively, the Laplace and minimal-surface equations.  The prescribed
mean-curvature equation is
\begin{equation}
\mathcal Mu=(1+|Du|^2)\Delta u-D_iuD_juD_{ij}u
=nH(1+|Du|^2)^{3/2};
\tag{10.7}
\end{equation}
here $\lambda=1$, $\Lambda=1+|p|^2$, and $\mathcal E=|p|^2$.
The subsonic potential-flow equation
\begin{equation}
\Delta u-
\frac{D_iuD_ju}{1-\frac{\gamma-1}{2}|Du|^2}D_{ij}u=0
\tag{10.8}
\end{equation}
is elliptic when $|Du|<(2/(\gamma+1))^{1/2}$.  The capillarity equation is
\begin{equation}
\operatorname{div}\frac{Du}{\sqrt{1+|Du|^2}}=\kappa u.
\tag{10.9}
\end{equation}

\section{Comparison Principle}

\begin{theorem}[Comparison Principle]
\label{gilbarg-trudinger.thm.10.1.comparison-principle}\dcref{ch10:10.1}\group{chapter-10}
\source{gilbarg-trudinger:page-0275}
\source{gilbarg-trudinger:page-0276}
Let $u,v\in C^0(\overline\Omega)\cap C^2(\Omega)$ satisfy
$Qu\ge Qv$ in $\Omega$ and $u\le v$ on $\partial\Omega$.  Assume:
\begin{enumerate}
\item[(i)] $Q$ is locally uniformly elliptic with respect to either $u$ or $v$;
\item[(ii)] $a^{ij}$ is independent of $z$;
\item[(iii)] $b(x,z,p)$ is nonincreasing in $z$ for fixed $(x,p)$;
\item[(iv)] $a^{ij}$ and $b$ are continuously differentiable in $p$.
\end{enumerate}
Then $u\le v$ in $\Omega$.  If $Qu>Qv$, the conclusion is $u<v$ and
condition \textup{(iv)} is unnecessary.
\end{theorem}
\begin{proof}
\sketch Suppose first that $Q$ is elliptic with respect to $u$ and put
$w=u-v$.  On $\{w>0\}$, the mean-value theorem and condition \textup{(iii)}
give a linear inequality
\[
a^{ij}(x,u,Du)D_{ij}w+b^iD_iw\ge0
\]
with locally bounded $b^i$.  The weak maximum principle gives $w\le0$.
If $Q$ is elliptic with respect to $v$, apply the minimum principle instead.
For the strict assertion, a nonnegative interior maximum contradicts
$Qu>Qv$ directly.
\end{proof}

\begin{theorem}[Dirichlet Uniqueness]
\label{gilbarg-trudinger.thm.10.2.dirichlet-uniqueness}\dcref{ch10:10.2}\group{chapter-10}
\source{gilbarg-trudinger:page-0276}
\uses{gilbarg-trudinger.thm.10.1.comparison-principle}
Under the hypotheses of Theorem~\ref{gilbarg-trudinger.thm.10.1.comparison-principle},
if $Qu=Qv$ in $\Omega$ and $u=v$ on $\partial\Omega$, then $u=v$ in
$\Omega$.
\end{theorem}
\begin{proof}
\sketch Apply the comparison principle in both directions.
\end{proof}

Theorems 10.1 and 10.2 also hold for
$u,v\in C^0(\overline\Omega)\cap C^1(\Omega)
\cap W^{2,n}_{\mathrm{loc}}(\Omega)$ by the strong-solution maximum principle.

\section{Maximum Principles}

\begin{theorem}[A Priori Bound]
\label{gilbarg-trudinger.thm.10.3.apriori-bound}\dcref{ch10:10.3}\group{chapter-10}
\source{gilbarg-trudinger:page-0277}
Let $Q$ be elliptic in $\Omega$ and suppose that $\mu_1,\mu_2\ge0$ satisfy
\begin{equation}
\frac{b(x,z,p)\operatorname{sign}z}{\mathcal E(x,z,p)}
\le\frac{\mu_1|p|+\mu_2}{|p|^2}
\tag{10.10}
\end{equation}
for $(x,z,p)\in\Omega\times\R\times\R^n$.  If
$u\in C^0(\overline\Omega)\cap C^2(\Omega)$ and $Qu\ge0$, then
\begin{equation}
\sup_\Omega u\le\sup_{\partial\Omega}u^++C\mu_2,
\qquad C=C(\mu_1,\operatorname{diam}\Omega).
\tag{10.11}
\end{equation}
For a solution, the corresponding estimate holds for $|u|$.
\end{theorem}
\begin{proof}
\sketch Place $\Omega$ in $0<x_1<d$ and compare $u$ on $\{u>0\}$ with
\[
v=\sup_{\partial\Omega}u^+
+\mu_2(e^{\alpha d}-e^{\alpha x_1}),\qquad \alpha\ge\mu_1+1.
\]
Condition \textup{(10.10)} gives $\bar Qv<0\le\bar Qu$ for the operator
with coefficients frozen along $u$.  Theorem~\ref{gilbarg-trudinger.thm.10.1.comparison-principle}
then gives \textup{(10.11)}; let $\mu_2\downarrow0$ for the limiting case.
\end{proof}

For uniformly elliptic operators, \textup{(10.10)} is equivalent to a bound
of the form
\begin{equation}
\frac{b(x,z,p)\operatorname{sign}z}{\lambda(x,z,p)}
\le\mu_1|p|+\mu_2.
\tag{10.12}
\end{equation}
For general elliptic operators put
$\mathscr D=\det[a^{ij}]$ and $\mathscr D^*=\mathscr D^{1/n}$.

\begin{theorem}[Aleksandrov-Type Bound]
\label{gilbarg-trudinger.thm.10.4.aleksandrov-type-bound}\dcref{ch10:10.4}\group{chapter-10}
\source{gilbarg-trudinger:page-0278}
Suppose $Q$ is elliptic and
\begin{equation}
\frac{b(x,z,p)\operatorname{sign}z}{\mathscr D^*}
\le\mu_1|p|+\mu_2
\tag{10.13}
\end{equation}
for constants $\mu_1,\mu_2\ge0$.  If
$u\in C^0(\overline\Omega)\cap C^2(\Omega)$ and $Qu\ge0$, then
\begin{equation}
\sup_\Omega u\le\sup_{\partial\Omega}u^++C\mu_2,
\qquad C=C(\mu_1,\operatorname{diam}\Omega).
\tag{10.14}
\end{equation}
For $Qu=0$, the same estimate holds with $u,u^+$ replaced by
$|u|,|u|$.
\end{theorem}
\begin{proof}
\sketch On $\{u>0\}$, condition \textup{(10.13)} turns $Qu\ge0$ into
the hypothesis of the Aleksandrov maximum principle.  Apply it to $u$ and,
for a solution, to $-u$.
\end{proof}

\begin{theorem}[Integral Bound]
\label{gilbarg-trudinger.thm.10.5.integral-bound}\dcref{ch10:10.5}\group{chapter-10}\source{gilbarg-trudinger:page-0278}
Let $\Omega$ be bounded and $Q$ elliptic.  Suppose nonnegative
$g\in L^n_{\mathrm{loc}}(\R^n)$ and $h\in L^n(\Omega)$ satisfy
\begin{equation}
\frac{b(x,z,p)\operatorname{sign}z}{n\mathscr D^*}
\le\frac{h(x)}{g(p)}
\tag{10.15}
\end{equation}
and
\begin{equation}
\int_\Omega h^n\,dx<\int_{\R^n}g^n\,dp=:g_\infty.
\tag{10.16}
\end{equation}
If $u\in C^0(\overline\Omega)\cap C^2(\Omega)$ and $Qu\ge0$, then
\begin{equation}
\sup_\Omega u\le\sup_{\partial\Omega}u^++C\operatorname{diam}\Omega.
\tag{10.17}
\end{equation}
For $Qu=0$, the analogous estimate holds for $|u|$.  If $g>0$ and
\[
G^{-1}(t)=\int_{B_t(0)}g^n\,dp,
\]
one may take $C=G(\int_\Omega h^n)$.  The condition \textup{(10.16)} is
void when $g_\infty=\infty$.
\end{theorem}
\begin{proof}
\sketch Apply the integral Aleksandrov estimate to $u$ on $\{u>0\}$;
the displayed inverse distribution formula gives the stated constant.  Apply
the same argument to $-u$ for a solution.
\end{proof}

\begin{corollary}[Prescribed Mean Curvature Bound]
\label{gilbarg-trudinger.cor.10.6.prescribed-mean-curvature-bound}\dcref{ch10:10.6}\group{chapter-10}
\source{gilbarg-trudinger:page-0279}
\uses{gilbarg-trudinger.thm.10.5.integral-bound}
Let $u\in C^0(\overline\Omega)\cap C^2(\Omega)$ solve
\textup{(10.7)} in a bounded domain.  If
\begin{equation}
H_0:=\int_\Omega|H(x)|^n\,dx<\omega_n,
\tag{10.18}
\end{equation}
then
\begin{equation}
\sup_\Omega|u|\le\sup_{\partial\Omega}|u|
+C(n,H_0)\operatorname{diam}\Omega.
\tag{10.19}
\end{equation}
\end{corollary}
\begin{proof}
\sketch For \textup{(10.7)},
$\mathscr D=(1+|p|^2)^{n-1}$.  Use
$g(p)=(1+|p|^2)^{-(n+2)/(2n)}$, for which
$\int_{\R^n}g^n=\omega_n$, in Theorem~\ref{gilbarg-trudinger.thm.10.5.integral-bound}.
\end{proof}

Theorems 10.3--10.5 and Corollary 10.6 remain valid for
$u\in C^0(\overline\Omega)\cap W^{2,n}_{\mathrm{loc}}(\Omega)$.

\section{Failure When the Principal Coefficients Depend on the Solution}

\begin{example}[Smooth nonuniqueness counterexample]
\label{gilbarg-trudinger.ex.10.3.solution-dependent-counterexample}\dcref{ch10:10.3}\group{chapter-10}\source{gilbarg-trudinger:page-0280}
On $\Omega=\{x\in\R^n:1<|x|<2\}$, consider
\begin{equation}
\mathcal Qu=\Delta u+g(r,u)\frac{x_ix_j}{r^2}D_{ij}u,
\qquad r=|x|.
\tag{10.20}
\end{equation}
Choose polynomials $v,w$ on $[1,2]$ such that
\[
v(1)=w(1),\quad v(2)=w(2),\quad v',w'>0,\quad v'',w''<0,
\]
\[
v'(1)<w'(1),\quad v'(2)>w'(2),\quad
\frac{v''}{v'}=\frac{w''}{w'}\text{ at }r=1,2.
\]
For $v(r)\le z\le w(r)$ define
\[
f(r,z)=\frac{z-v}{w-v}
\left(\frac{v''}{v'}-\frac{w''}{w'}\right)-\frac{v''}{v'},
\qquad
g(r,z)=-1+\frac{n-1}{rf(r,z)}.
\]
After a smooth extension in $z$, $\mathcal Q$ is uniformly elliptic with
smooth coefficients, while the distinct radial functions $v(|x|)$ and
$w(|x|)$ solve $\mathcal Qv=\mathcal Qw=0$ and have the same boundary data.
\end{example}

\section{Comparison for Divergence-Form Operators}

A weakly differentiable $u$ satisfies $Qu\ge0$ for \textup{(10.5)} when
\begin{equation}
Q(u,\varphi):=\int_\Omega
\{A(x,u,Du)\cdot D\varphi-B(x,u,Du)\varphi\}\,dx\le0
\tag{10.21}
\end{equation}
for every nonnegative $\varphi\in C_0^1(\Omega)$.

\begin{theorem}[Comparison for Divergence Form]
\label{gilbarg-trudinger.thm.10.7.comparison-divergence-form}\dcref{ch10:10.7}\group{chapter-10}
\source{gilbarg-trudinger:page-0281}
Let $u,v\in C^1(\overline\Omega)$ satisfy $Qu\ge0$, $Qv\le0$ in
$\Omega$ and $u\le v$ on $\partial\Omega$.  Suppose $A,B$ are
continuously differentiable in $(z,p)$, $Q$ is elliptic, and $B$ is
nonincreasing in $z$.  Then $u\le v$ if any one of the following holds:
\begin{enumerate}
\item[(i)] $A$ is independent of $z$;
\item[(ii)] $B$ is independent of $p$;
\item[(iii)] the matrix
\[
\begin{bmatrix}
D_{p_j}A^i&-D_{p_j}B\\
D_zA^i&-D_zB
\end{bmatrix}
\]
is nonnegative on $\Omega\times\R\times\R^n$.
\end{enumerate}
\end{theorem}
\begin{proof}
\sketch Put $w=u-v$, $u_t=tu+(1-t)v$, and average the derivatives of
$A,B$ along $u_t$.  For every nonnegative test function,
\begin{equation}
\begin{aligned}
0\ge Q(u,\varphi)-Q(v,\varphi)
=\int_\Omega\{&(a^{ij}D_jw+b^iw)D_i\varphi\\
&-(c^iD_iw+dw)\varphi\}\,dx.
\end{aligned}
\tag{10.22}
\end{equation}
Thus $w$ is a subsolution of
$D_i(a^{ij}D_jw+b^iw)+c^iD_iw+dw$.  In case \textup{(i)}, the weak
maximum principle applies because $b^i=0$.  In case \textup{(ii)}, test by
$w^+/(w^++\varepsilon)$; Poincare's inequality bounds
$\log(1+w^+/\varepsilon)$ uniformly, so $w^+=0$ as
$\varepsilon\downarrow0$.  In case \textup{(iii)}, the matrix inequality
and the test $w^+$ give $|Dw^+|\le Cw^+$; applying this to the same
logarithm again forces $w^+=0$.
\end{proof}

Under condition \textup{(i)}, it is enough that
$u,v\in C^0(\overline\Omega)\cap C^1(\Omega)$ and that the coefficient
derivatives are continuous; analogous extensions hold in the other cases
under uniform structure bounds.

\section{Maximum Principles for Divergence-Form Operators}

Assume that, for some $\alpha\ge1$ and nonnegative constants
$a_1,a_2,b_0,b_1,b_2$,
\begin{equation}
\begin{aligned}
p\cdot A(x,z,p)&\ge|p|^\alpha-|a_1z|^\alpha-a_2^\alpha,\\
B(x,z,p)\operatorname{sign}z&\le
\begin{cases}
b_0|p|^{\alpha-1}+|b_1z|^{\alpha-1}+b_2^{\alpha-1},&\alpha>1,\\
b_0,&\alpha=1.
\end{cases}
\end{aligned}
\tag{10.23}
\end{equation}

\begin{lemma}[Divergence Form Maximum Bound]
\label{gilbarg-trudinger.lem.10.8.divergence-form-maximum-bound}\dcref{ch10:10.8}\group{chapter-10}
\source{gilbarg-trudinger:page-0283}
Let $u\in C^0(\overline\Omega)\cap C^1(\Omega)$ satisfy $Qu\ge0$ and
assume \textup{(10.23)}.  Then
\begin{equation}
\sup_\Omega u\le C\{\|u^+\|_\alpha
+(a_1+b_1)\sup_{\partial\Omega}u^++a_2+b_2\}
+\sup_{\partial\Omega}u^+,
\tag{10.24}
\end{equation}
where $C=C(n,\alpha,a_1,b_0,b_1,|\Omega|)$.
\end{lemma}
\begin{proof}
\sketch First take $u\le0$ on the boundary and set
$k=a_2+b_2$, $w=u^++k$, and
$\bar b=a_1^\alpha+b_0^\alpha+b_1^{\alpha-1}+1$.  Young's inequality
gives
\begin{equation}
p\cdot A\ge|p|^\alpha-\bar b|z+k|^\alpha,
\tag{10.25}
\end{equation}
together with the analogous bound for $(z+k)B\operatorname{sign}z$.
Testing \textup{(10.21)} by $w^\beta-k^\beta$ and applying Sobolev gives
an iteration
\[
\|w\|_{rs}\le(Cr)^{1/r}\bar b^{1/(\alpha r)}\|w\|_{r\alpha},
\quad r\ge1,
\]
which yields \textup{(10.24)}.  Subtract the positive boundary supremum
and exhaust $\Omega$ by interior domains to remove the initial assumptions.
\end{proof}

\begin{theorem}[Divergence Form A Priori Estimate]
\label{gilbarg-trudinger.thm.10.9.divergence-form-apriori-estimate}\dcref{ch10:10.9}\group{chapter-10}
\source{gilbarg-trudinger:page-0283}
\uses{gilbarg-trudinger.lem.10.8.divergence-form-maximum-bound}
Assume \textup{(10.23)} with $\alpha>1$, $b_1=0$, and either $b_0=0$
or $a_1=0$.  If $u\in C^0(\overline\Omega)\cap C^1(\Omega)$ and
$Qu\ge0$, then
\begin{equation}
\sup_\Omega u\le C\{a_2+b_2+a_1\sup_{\partial\Omega}u^+\}
+\sup_{\partial\Omega}u^+.
\tag{10.26}
\end{equation}
For $Qu=0$, the same estimate holds for $|u|$, with boundary term
$\sup_{\partial\Omega}|u|$.  Here
$C=C(n,\alpha,a_1,b_0,|\Omega|)$.
\end{theorem}
\begin{proof}
\sketch Reduce as in Lemma~\ref{gilbarg-trudinger.lem.10.8.divergence-form-maximum-bound}
to zero boundary data and $k=a_2+b_2>0$.  If $b_0=0$, test by
$k^{1-\alpha}-w^{1-\alpha}$ to bound
$\|D\log(w/k)\|_\alpha$ and hence $\sup w/k$.  If $a_1=0$, use
the reciprocal-power test based on $M-w+k$, obtaining
\begin{equation}
\int_\Omega\left|\log\frac{M}{M-w+k}\right|^\alpha dx\le C.
\tag{10.27}
\end{equation}
The logarithm satisfies the hypotheses of Lemma~\ref{gilbarg-trudinger.lem.10.8.divergence-form-maximum-bound},
which again gives $M\le Ck$.  Let $k\downarrow0$, then restore the boundary
values and use $-u$ for solutions.
\end{proof}

\begin{theorem}[Small-Data Divergence Form Bound]
\label{gilbarg-trudinger.thm.10.10.small-data-divergence-form-bound}\dcref{ch10:10.10}\group{chapter-10}
\source{gilbarg-trudinger:page-0287}
Assume \textup{(10.23)} and let
$u\in C^0(\overline\Omega)\cap C^1(\Omega)$ satisfy $Qu\ge0$.  There is
$C_0=C_0(\alpha,n)>0$ such that, if
\begin{equation}
(a_1^\alpha+b_0^\alpha+b_1^{\alpha-1})|\Omega|^{\alpha/n}<C_0,
\tag{10.28}
\end{equation}
then
\begin{equation}
\sup_\Omega u^+\le
C\{(a_1+b_1)\sup_{\partial\Omega}u^++a_2+b_2\}
+\sup_{\partial\Omega}u^+,
\tag{10.29}
\end{equation}
where $C=C(n,\alpha,a_1,b_0,b_1,|\Omega|)$.  For a solution, the same
bound holds with $u^+$ replaced by $|u|$.
\end{theorem}
\begin{proof}
\sketch By Lemma~\ref{gilbarg-trudinger.lem.10.8.divergence-form-maximum-bound},
it suffices to control $\|u^+\|_\alpha$.  With zero boundary data, test
\textup{(10.21)} by $u^+$.  Young's and Poincare's inequalities bound the
$L^\alpha$ norm by itself times the left side of \textup{(10.28)}, plus
the $a_2,b_2$ terms.  The smallness hypothesis absorbs the first term;
Lemma 10.8 then gives \textup{(10.29)}.
\end{proof}

When $\alpha=1$, $b_1,b_2$ are absent.  The sharp inequality
\begin{equation}
\int_\Omega|v|\,dx\le\frac1n
\left(\frac{|\Omega|}{\omega_n}\right)^{1/n}
\int_\Omega|Dv|\,dx,
\qquad v\in W^{1,1}_0(\Omega),
\tag{10.30}
\end{equation}
allows $C_0(1,n)=n\omega_n^{1/n}$.  In particular, the divergence form
of prescribed mean curvature is
\begin{equation}
\operatorname{div}\frac{Du}{\sqrt{1+|Du|^2}}=nH.
\tag{10.31}
\end{equation}
If
\begin{equation}
H_0=\sup_\Omega|H|<\left(\frac{\omega_n}{|\Omega|}\right)^{1/n},
\tag{10.32}
\end{equation}
then every subsolution satisfies, and every solution satisfies in absolute
value,
\begin{equation}
\sup_\Omega u\le\sup_{\partial\Omega}u+C(n,|\Omega|,H_0).
\tag{10.33}
\end{equation}

The structure coefficients may instead be nonnegative measurable functions.
If $a_1,a_2,b_0,b_1,b_2\in L^q(\Omega)$, $q\ge\alpha$ and $q>n$, the
preceding estimates hold with the corresponding $L^q$ norms.  The smallness
condition becomes
\begin{equation}
\|a_1^\alpha+b_0^\alpha+b_1^{\alpha-1}\|_\beta<C_0,
\qquad \beta=\max(1,n/\alpha).
\tag{10.34}
\end{equation}
For \textup{(10.31)}, the condition
\begin{equation}
\int_\Omega|H|^n\,dx<\omega_n
\tag{10.35}
\end{equation}
implies \textup{(10.33)} with $H_0=\|H\|_n$.  These divergence-form
results also hold for $u\in W^{1,1}_*(\Omega)$ whenever the displayed
quantities are defined.
